%=====================================================================
%  Chapter 8 — Simulation, Testing and Validation
%=====================================================================
\chapter{Simulation, Testing and Validation}
\label{chap:validation}

\section{Introduction}
This chapter evaluates the complete system. Because the robot is driven by
stiff position servos and a fall in hardware risks the 3D-printed
structure, validation is carried out primarily on the Gazebo digital twin
introduced in Chapter~\ref{chap:software}, where behaviour can be observed
repeatably and safely. The evaluation proceeds bottom-up, mirroring the
order in which the system was built: first standing stability, then active
balance, then the walking gait, and finally the natural-language command
layer.

\section{Simulation Environment and Method}
The tests run in Gazebo Harmonic with the DART physics engine simulating
gravity, contacts and joint dynamics. The simulator can be launched with a
graphical view for interactive testing --- where the robot can be reset and
pushed to probe disturbance rejection --- or headless for faster, repeatable
runs. The principal quantitative measure throughout is the \textbf{torso
tilt}: the pitch and roll estimated by the complementary filter from the
inertial sensor (Chapter~\ref{chap:kinematics}). A tilt that stays small
indicates an upright, stable robot; a tilt that grows without bound
indicates a fall. This single signal is used to judge standing, balancing
and walking alike.

\section{Standing Stability}
The first test confirms the foundation: with no controller active, the bare
robot is spawned in its geometric standing pose. It stands stably, with the
torso tilt holding at roughly half a degree --- confirming that the
kinematic model, the mass and inertia properties, and the foot contact are
all sound, and that any later instability comes from the motion, not from a
broken model. Running the full software chain with the gait and balance
controllers disabled likewise leaves the robot standing, establishing that
the end-to-end pipeline is correct before any active control is added.

\section{Balance Validation}
\label{sec:val-balance}
The second test evaluates active balance on a more demanding pose: a deep
knee-bend crouch, which --- unlike the rigid straight-leg stand --- is an
inverted pendulum and is unstable without feedback. Two ingredients make it
hold. First, the target pose is \emph{eased into} over about a second
rather than commanded instantly, because the simulated position controller
otherwise snaps the joints and the robot lurches before the balancer can
react. Second, the LQR ankle controller of Chapter~\ref{chap:kinematics} is
enabled. With both, the robot holds the crouch stably at around four
degrees of tilt for twenty seconds and more --- against a fall to tens of
degrees within two seconds when balance is off. The controller is
effective at holding a settled stance and rejecting gentle pushes; its
authority is bounded by the ankle correction limit, so it cannot catch a
violent lurch, only maintain a stance that is already near equilibrium.

\begin{figure}[ht]
  \centering
  \includegraphics[width=0.85\textwidth]{fig8.1.png}
  \caption{The robot standing in the Gazebo simulation, used as the baseline
           for the balance and gait tests.}
  \label{fig:val-stand}
\end{figure}

\section{Gait Validation}
\label{sec:val-gait}
The third test exercises the full walking pipeline of
Chapter~\ref{chap:kinematics}. The pipeline behaves correctly in its parts:
it produces a valid ZMP reference, a smooth centre-of-mass trajectory, and
properly alternating foot motions in which a foot only lifts while the
centre of mass is over the opposite, supporting foot.

In terms of whole-body behaviour, the honest result is a partial success.
On good runs the robot \textbf{marches in place} for several balanced steps
--- on the order of twelve to sixteen seconds --- holding a torso tilt in the
range of three to seven degrees, and the capture-point step adjustment
demonstrably catches near-falls, recovering from tilt spikes that would
otherwise end the run. However, the march is \emph{marginally} stable and
varies from run to run on identical code: a slow lateral oscillation tends
to build up and eventually tips the robot, and enabling true forward motion
tips it sooner. This is not a tuning oversight but a method-level limit:
stiff position servos without torque sensing, combined with only marginal
lateral balance during single support, place robust dynamic walking beyond
the reach of the present ankle-and-hip strategy. Closing the gap calls for
full capture-point footstep control or a learned policy, identified as
future work in Chapter~\ref{chap:conclusion}.

\section{Command-Layer Validation}
\label{sec:val-ai}
The natural-language command layer of Chapter~\ref{chap:ai} was validated
end-to-end through the live local model. Plain-English requests are
correctly grounded into the right command and executed: waving, bowing,
sitting and returning to a standing pose all fire and report back, and the
robot's spoken-style confirmation matches the action performed. The
supporting mechanisms were verified individually:
\begin{itemize}[itemsep=2pt,topsep=4pt]
  \item the \textbf{emergency stop} interrupts a motion already in progress
        --- a walk in flight is cancelled and the robot snaps to a safe pose
        --- confirming that the separate-context stop path is a genuine
        interrupt;
  \item the \textbf{safety validator} was unit-tested across its cases (a
        level walk is allowed, an immediate second motion is rate-limited,
        a walk on a tilted base is refused, and \texttt{stop} and
        \texttt{get\_status} are never blocked);
  \item the \textbf{feedback loop} was verified with a scripted scenario in
        which a refused command leads the model to steady the robot and then
        retry successfully, rather than give up.
\end{itemize}
A noteworthy bug was found and fixed during this validation: under
continuous status polling, the stop/status path and the dispatcher
initially shared one ROS\,2 context and corrupted each other's internal
state. Giving the stop/status path its own context resolved it --- an
example of the concurrency care that a real interactive robot interface
requires. The main residual limitation is the small (3-billion-parameter)
model itself, which occasionally mis-formats a tool call or omits an
argument; this is a known limitation of small models rather than of the
architecture.

\section{Results Summary}
Table~\ref{tab:results} summarises the validation outcomes.

\begin{table}[ht]
  \centering
  \caption{Summary of validation results.}
  \label{tab:results}
  \small
  \renewcommand{\arraystretch}{1.3}
  \begin{tabularx}{\textwidth}{@{}lXl@{}}
    \toprule
    \textbf{Capability} & \textbf{Method} & \textbf{Result} \\
    \midrule
    Standing stability & Bare robot, no control & Stable ($\sim$0.5\textdegree{} tilt) \\
    Static crouch      & Closed-form IK + LQR balance + ease-in & Held 20\,s+ ($\sim$4\textdegree) \\
    March in place     & ZMP preview gait + balance + capture point & 12--16\,s, 3--7\textdegree{} (good runs) \\
    Forward walking    & Gait with forward speed & A few steps, then tips --- not robust \\
    NL commands        & Local LLM tool-use $\to$ action & Grounded and executed correctly \\
    Emergency stop     & Service on isolated context & Interrupts in-flight motion \\
    Safety validator   & Tilt gate + rate limit & Verified (unit + live) \\
    Feedback loop      & Agentic recover-then-retry & Verified (scripted) \\
    \bottomrule
  \end{tabularx}
\end{table}

\section{Conclusion}
The validation confirms a system that is sound from the ground up: a correct
physical model, an effective standing balance, a functional gait pipeline
that marches in place, and a natural-language command layer that works
end-to-end through a local model with safety gating. It also delineates,
honestly, the project's frontier: robust dynamic forward walking on stiff
hobby servos remains unsolved and is the natural target of future work. The
next chapter draws the general conclusions of the project and lays out
those perspectives.
